% Fresh exhaustive canonical-source assembly, 2026-10-22.
\section{The complete canonical local certificate and global source dispatch}
\label{sec:complete-source-routing}

\paragraph{Quantifiers and the full-parent target.}
Let \(K_c,K_t\) be arbitrarily chosen, separately labelled real closed
extensions of \(\mathbb R\).  There is no algebraic-disjointness
assumption: the two factors may be labelled copies of the same field.
Put \(T=K_c\otimes_{\mathbb R}K_t\).
Consider a physical canonical source over \(K_c\) with an actual
finite noncommuting stationary zero and one of the exact patterns
\[
\begin{array}{ll}
{\rm C1}:&d,m,C,A,\alpha>0,\quad\lambda=n=0,\quad
             D,\beta\geq0,\quad D+\beta>0,\\
{\rm C2}:&d,m,C,A,D,\alpha>0,\quad\lambda=\beta=0,\quad n\geq0.
\end{array}\tag{LC1}
\]
These are precisely the two elementary source hypotheses produced
by the actual extreme-source reduction SE15.  For finite independent
testing half-angles and any state \((W,X,Y,Z)\in K_t^4\), the complete
parent is
\[
\begin{split}
 F={}&4\lambda D_u(W+vX)^2+4nD_uD_vX^2
       +4m(vW-X+uvY-uZ)^2\\
 &+4CD_v(uW-Y)^2+4AD_uD_vY^2
       +(4D+2d)(uW+uvX-Y-vZ)^2\\
 &+(4\alpha+2d)D_u(vY-Z)^2
       +(4\beta+2d)D_v(X+uZ)^2-2dD_uD_vZ^2 .
\end{split}\tag{LC2}
\]
Every coefficient times a testing expression denotes its separated
tensor product.  The homogeneous angle form of LC2, rather than
an affine limit obtained by a cancelled zero coordinate, defines
the literal projective endpoints.

\begin{theorem}[Complete local certificate]
For every source in LC1 and every independent projective testing pair
and real or complex testing state, the entire homogeneous parent is
a finite sum of squares in \(T\).
This includes every coefficient boundary permitted in LC1, every
positive source opening, all relative source and testing rates, and
all state ratios.  The source contact is never identified with the
testing point.
\end{theorem}

\subsection{Normalization and actual native entry}

Divide by the positive pure-source \(d\).  This operation is reversible
in a tensor sum of squares by multiplication of each root by
\(\sqrt d\otimes1\).  Before or after this normalization the invariant
total-scale quantity is
\[
 \delta_{\rm tot}
  =\frac d{d+\lambda+n+m+C+A+D+\alpha+\beta}.
\tag{LC3}
\]
It is positive and at most one.  Its positive-unit and positive
infinitesimal alternatives are exhaustive.  In the first alternative
the complete positive-scale theorem AU1--AU19 applies after EN.
The second alternative retains the hypothesis
\(\delta_{\rm tot}\in\mathrm{Inf}_{c,>0}\) in every residual theorem
used below.  Setting \(d=1\) does not erase that hypothesis.

The theorem EN1--EN16 derives the full native signs and exact rational
source identification from physicality and the actual finite contact.
Its proof uses SE's coercive Laurent support, strict single-field
contact Hessian, contact sufficiency and event-row deletion implication.
It does not assume any tensor certificate being assembled here.
For C1, after the exact Alice source reflection, its native coordinates
satisfy
\[
\begin{gathered}
 0<k\leq1,\quad \eta,\zeta>0,\quad
 r=1-a\eta\in(0,1),\quad \theta=t\zeta\in(0,1),\\
 0<a,t<2,\quad x=k^2,\quad h=1-\theta,\quad
 E=1-x\eta h>0,\quad Q=1-x\zeta r>0,\\
 M=1+xr-t>0,\quad J=1+xh-a>0,\qquad
 e=k\sqrt{\eta/\zeta}>0 .
\end{gathered}\tag{LC4}
\]
The original \(D,\beta\) remain nonnegative, and either may be zero.
For C2, EN gives exactly CT2 and its unique actual stationary
coefficient formulas, including \(n=0\).
All these implications hold over every \(K_c\) by the first-order
transfer proved in EN; no topological compactness of a
non-Archimedean field is being presumed.

\subsection{Every weak-total C1 source}

Use the native parameters LC4, the actual original coefficients, and
\(\delta_{\rm tot}\) positive infinitesimal.
Here I, U and O mean a positive infinitesimal, positive bounded unit,
and positive unbounded source quantity, respectively.
Since \(0<k\leq1\), \(k\) has type I or U.
The following disjoint scale dispatch is exhaustive:
\[
\begin{array}{c|c|c|l}
k&\zeta&\text{remaining source type}&\text{actual theorem}\\ \hline
{\rm U}&\text{any}&\eta,\zeta:{\rm I}&{\rm AS}\\
{\rm U}&\text{any}&\eta,\zeta\text{ not both I}&
                      \text{ordinary-Alice join below}\\
{\rm I}&{\rm U\ or\ O}&\eta>0&
                      {\rm JC}\text{ with bare source scope}\\
{\rm I}&{\rm I}&\eta:{\rm I}&
                      \text{uniform all-}e\text{ join}\\
{\rm I}&{\rm I}&\eta:{\rm O}&{\rm C1.1}\\
{\rm I}&{\rm I}&\eta:{\rm U},\ e:{\rm I}&{\rm C1.3}\\
{\rm I}&{\rm I}&\eta,e,r:{\rm U}&{\rm BE}\\
{\rm I}&{\rm I}&\eta,e:{\rm U},\ r:{\rm I}&{\rm CA}\\
{\rm I}&{\rm I}&\eta,r:{\rm U},\ e:{\rm O}&{\rm UE}_{\rm old}\\
{\rm I}&{\rm I}&\eta:{\rm U},\ e:{\rm O},\ r:{\rm I}
                                                  &{\rm C1.2}.
\end{array}\tag{LC5}
\]
The rows using \(r\) are exhaustive because \(0<r<1\).
The positive number \(e\) has exactly the three types shown.
No comparison with an independent testing quantity occurs in LC5.

\paragraph{Ordinary Alice: the actual UI, UU and UH split.}
Suppose \(k\) is a unit and \(\eta,\zeta\) are not both infinitesimal.
If \(\eta\) is infinitesimal, then \(r\sim1\) and \(Q>0\)
make \(\zeta\) bounded.  Thus \(\zeta\) is a unit in this case.
If \(\zeta\) is infinitesimal, then \(\theta\) is infinitesimal,
\(h\sim1\), and \(E>0\) make \(\eta\) bounded; hence \(\eta\)
is a unit.  In particular the two possibilities
\[
 (\eta:{\rm I},\zeta:{\rm O}),\qquad
 (\eta:{\rm O},\zeta:{\rm I})
\tag{LC6}
\]
are impossible when \(k\) is a unit.
Since \(v_c^2=x\eta\zeta\), the remaining possibilities are:
one of \(\eta,\zeta\) infinitesimal and the other a unit, giving
actual UI; both units, giving actual UU; or both noninfinitesimal
with at least one unbounded, giving actual UH.

The actual UI source theorem and its exact native stabilizer cover
the first possibility.  More explicitly, the exact AS3--AS5
full-parent reflection reduces it to \(\eta\) a unit and
\(\zeta\) infinitesimal.  This reflection preserves the C1 omission
pair and the original nonnegative coefficients; its new native
parameters and gaps are checked in that identity.  Its Alice
parameter remains a unit, but a second assertion \(k\leq1\) is
not imposed on the reflected representative.
Use the complete CU1--CU13 theorem in
\emph{ui\_all\_tests.tex}, together with CR1--CR13 in the actual
complementary-testing body.  For these displayed CU source
quotients rescale by the positive pure-source factor to its
normalization \(C=2\); the native parameters are unchanged,
and this local normalization is restored afterwards.
In CU put \(g^2=\zeta\).
The exact source quotients make \(C/a,A/a\) strong,
\(d/a,\alpha/(ag^2)\) units, and \(m/a,\beta/a\) large.
CU2 retains the entire paired \(m,\beta\) normal packets;
CU3 retains the centered \(C,A,D\) excesses.
CU5 is the whole remaining quadratic state polynomial on
its bounded six-root frame.  CU6--CU9 explicitly divide all
complete first-normal coefficients and the scalar correction in the
contact ideals.  Their exact tensor zero redistributes the
two angular pivots; CU10 computes the full matched state
Schur Hessian.  Its positive determinant uses the original
\(D\geq0\) numerator.  The possible divergent \(A,D\) pair has
independent hard rows, so its constrained finite-cap limit
keeps the angular components of those rows.
At unmatched Alice residue the normalized Laurent coefficients
retain coercivity and strict convexity.  The actual real
two-block lemma UA1--UA3 is applied using the exact relation
\(L(gZ)-Q(vZ)=0\) between its root products, not by inverting
a determinant in \(T\).  CR supplies all complementary tests,
especially its CR5a five-root reciprocal-Bob frame retaining
the weak \(\alpha\) packet.  This proves all UI tests and states.
Literal \(D=0\) is retained; \(\beta=0\) is impossible here by
the displayed lower bound on the actual \(\beta/a\).
For actual UU use the full NC1--NC7 source proof in
\emph{c1\_source\_formulas.tex}.
It first proves the uniform native gap \(E\geq1/18\).
The nonnegative \(D,\beta\) formulas and their exact elimination
identity give
\[
 C,A,m,\alpha\text{ strong},\qquad
                    D\text{ or }\beta\text{ strong}.
\tag{LC7}
\]
In more detail \(x,v_c^2,\eta,\zeta\) are units; the specialization
\(B_\beta|_{\theta=1}=-1-x\zeta(1-r)\) makes \(h\) a unit.
The original numerator polynomials \(P_4,P_\alpha,P_d\) are
positive units.  If both \(D,\beta\) were weak, the left side of
the exact NC7 elimination would be infinitesimal and its
explicit positive right side a unit.  This is impossible.
The five consecutive strong events in LC7 therefore give the
actual five-event all-tests proof, with its detailed A6
noncommuting testing complement.
This is a source theorem at all upper event scales, not an
application of a bounded-coefficient theorem to unbounded data.
For actual UH the native NI8--NI10 implications give two
alternatives: \(\eta\) unbounded with \(\zeta\) a unit, or
\(\zeta\) unbounded with \(\eta\) noninfinitesimal.
The exact \(\beta\)-numerator identity proves \(\eta h\) a
unit in both.  Thus, with \(\sigma=1/\eta\), \(z=1/\zeta\)
and \(H=h/\sigma\), the source is respectively
\(\sigma\) infinitesimal and \(z\) a unit, or \(z\)
infinitesimal and \(\sigma\) bounded, always with \(H\)
a unit.  These are precisely the printed UH1 alternatives;
no bound on the ratio of two source infinitesimals is added.
The actual four-consecutive-event theorem CE and its
matched weak outer frames CE17--CE21 remove the remote tests.
The remaining strip has unit testing \(u\) and infinitesimal
\(w=1/v\).  The complete homogeneous biquadratic conductor
\emph{uh\_all\_tests.tex} applies there.
With \(b=r/z\), use its scaled weights
\(\widehat m=m/z\), \(\widehat C=C/z\),
\(\widehat A=Az\) and \(\widehat\alpha=\alpha z\).
Their products
\(\widehat mM,\widehat C(1-r),\widehat A b,\widehat\alpha Q\)
are bounded and cancel the opening poles before the full
scalar and first-normal collection.
At a matched Alice residue its exact angular zero and
the full hard Schur calculation redistribute the positive
angular pivots.  At an unmatched residue its raw angular
product relation gives the finite real two-block allocation.
The single and double opening limits retain all divergent
centered rows; their independent hard-row determinants and
bounded hard minimizers justify the constrained limits.
This is a full-parent construction, including the
\(J=\theta=0\) residue face.
Both \(D=0\) and \(\beta=0\) are retained where permitted by the
actual source equations.  Thus no assertion that UI or UH alone
already covers UU is needed.

\paragraph{The uniform source and the four earlier exits.}
In the fourth row of LC5, \(k,\sqrt\eta,\sqrt\zeta\) are all
infinitesimal, while \(e=k\sqrt{\eta/\zeta}\) is arbitrary.
The new uniform all-\(e\) join UE1--UE4 applies exactly.
For \(e\) infinitesimal it uses the fresh complete UJ proof,
with its full bounded-ratio and inverse-ratio state charts.
Otherwise \(1/e\) is bounded, and the actual complete
bounded-Delta proof supplies its restricted monic module,
weighted full scalar, both entire first normals, hard block,
and complementary state chart.  Its outer current formulas
are joined in UE without an unprinted restriction on \(e\).
Neither a historical whole-uniform assertion nor a missing
old conductor label is a premise.

The AS, BE, CA and \({\rm UE}_{\rm old}\) rows use their actual
full-parent proofs, respectively
\path{arbitrary_state_ratio.tex},
\path{bounded_e_all_tests.tex},
\path{centered_a_unit_e.tex} and
\path{unit_eta_unbounded_e.tex}.
Their ordinary-angle, centered-A, contact-conductor and
outer-current subordinate bodies are included explicitly.
AS proves the needed strict original event signs on its source
domain.  BE and CA include \(\beta=0\).
The last proof derives \(\beta/d\) unbounded, so its intersection
with literal \(\beta=0\) is empty by an actual source estimate.
These proofs use the bare LC4 equations and their stated source
scales; no estimate comes from exclusion of an unnamed earlier
admission.

\paragraph{The remaining C1 rows and the complete C1.4 union.}
For C1.1 the actual source reduction, including the printed
\(k^2\eta\leq1/2+\mathrm{Inf}_c\) derivation, splits into
the unit-\(q\) and infinitesimal-\(q\) branches.  The complete
unit-\(q\) body and the protected-normal local and remote
infinitesimal-\(q\) bodies prove every testing/state chart.
C1.2 uses the actual centered-A bounded-Delta theorem, including
its full contact conductor and centered-A outer proof.
C1.3 uses the actual small-\(e\) all-tests theorem and its
explicit outer-deficit calculation.  These are full bare-source
statements with the scales printed in LC5.

For the third row of LC5, the source is exactly the native
C1.4 input \(k\) infinitesimal, \(\zeta\) noninfinitesimal,
with all original gaps and \(\delta_{\rm tot}\) infinitesimal.
We now give its bare-source join explicitly; this also explains
why a blanket earlier-complement predicate in the entered JC
statement is unnecessary for this invocation.

Put \(L=x\zeta\).  If \(L\) is unbounded, \(Q=1-Lr>0\)
forces \(r\) infinitesimal.  Since \(1-r=a\eta\) and
\(0<a<2\), \(\eta\) cannot be infinitesimal.
Unit \(\eta\) is the complete \emph{large\_l\_unit.tex}
exit.  For unbounded \(\eta\), put
\(\sigma=1/\eta\), \(H=h/\sigma\) and
\(M_\beta=\beta\theta/\sigma^2\).
Bounded \(H,M_\beta\), including \(\beta=0\), is the actual
JR/JB theorem.  Bounded \(H\) with unbounded \(M_\beta\)
is the full C14.1a SF conductor.  Unbounded \(H\) is
the C14.1b source: SF first treats \(\operatorname{res}h<1/2\);
US bounds \(k/\sigma\) and treats its unit case.
The BS/BC/GA and CE testing and state exits then leave
exactly the small-sphere source and the infinitesimal
testing chart stated in JC5.  The fresh AR theorem restores
the angular budget and all SP6 first normals there.
AR requires no bound on \(qu/w\), \(u/w\), \(W/Z\) or
\(uW/(wZ)\).  This closes the large-\(L\) alternative
without a comparison of independent source and test rates.

If \(L\) is a unit, the three types of \(\eta\) give,
respectively, the complete \emph{lunit\_eta\_small.tex}
theorem, the LU theorem in \emph{eta\_unit.tex},
and C14.2.  The first two use the actual LII parent
dictionary with their own native coefficient checks.
They do not identify the source with an arbitrary C2
source satisfying a narrower theorem.
C14.2 uses its explicit NC/LC source bounds,
complete \(C\)-opening extraction, HG/CT radial matrix,
sphere-weighted budgets, CS/CF/HG state joins and LA/LB
remote tests.  Every part has the full parent and all
normal terms.

If \(L\) is infinitesimal and \(\eta\) unbounded, then
\(\sigma=1/\eta\) is infinitesimal and this is C14.4.
Its source-only parameters and exact identity are
\[
\begin{gathered}
 S=\sigma+L,\quad \kappa=x/S,\quad
 \nu=\sigma/S,\quad \ell=L/S,\quad \nu+\ell=1,\quad
 \zeta=\ell/\kappa,\\
 Q_s=1-\ell Sr,\qquad
 \kappa Q_s+\ell M+\nu E=1,\qquad E\geq1/18.
\end{gathered}\tag{LC4a}
\]
Consequently \(\kappa\) is bounded.  The fractions
\(\nu,\ell\) cannot both be infinitesimal; if \(\kappa\)
is a unit, then \(\ell=\kappa\zeta\) is a unit.
Thus the five exact source cells are
\[
\begin{array}{c|c|c|l}
\kappa&\nu&\ell&\text{whole-parent theorem}\\ \hline
{\rm U}&{\rm U}&{\rm U}&{\rm FU}\text{ with its UKR matched chart}\\
{\rm U}&{\rm I}&{\rm U}&{\rm NU}\\
{\rm I}&{\rm U}&{\rm U}&{\rm FM}\\
{\rm I}&{\rm U}&{\rm I}&{\rm EL}\\
{\rm I}&{\rm I}&{\rm U}&{\rm SN,\ BN,\ UN}.
\end{array}\tag{LC4b}
\]
For the last cell put \(q_\nu=\nu/\ell\),
\(\tau=h/q_\nu\) and \(n_s=\beta/q_\nu^2\).
Bounded \(n_s\) gives SN and \(\tau\sim1\);
unbounded \(n_s\) with bounded \(\tau\) gives BN;
unbounded \(\tau\) gives UN, including infinitesimal
\(\theta\).  Each source is covered.
All these proofs start from the exact FM dictionary,
NC5 and their actual source/sphere equations.
Their SE/EB/CE testing exits have those same printed
source inequalities.  They retain all normal terms,
literal coefficient zeros where permitted, and every
source opening.

It remains to take \(L\) infinitesimal and \(\eta\)
bounded.  Unbounded \(\zeta\) is C14.3a.
Unit \(\zeta\) with infinitesimal \(\eta\) is C14.3b.
Their complete JF parent is combined with DG/LD normal
conductors and the actual completed \(j\) conductor;
the latter removes the earlier \(x/\eta\) restriction.
The old restricted follow-up assertion is not invoked.
DS/LS and SB supply all states and the remote tests.
The small-\(\eta\), bounded-\(a\) input is used only
through its proved angular paragraph, not its old
proposed radial continuation.

The remaining \(\eta,\zeta\) are units.  Use the
actual positive source gap
\[
 j=\eta(1+xh)-1+r=\eta J>0,\qquad
 j_0=\eta+r-1,\qquad j-j_0=\eta xh\in\mathrm{Inf}_c .
\tag{LC4c}
\]
It is bounded.  Unit \(j\) gives positive unit \(j_0\)
and the RO theorem.  Infinitesimal \(j\) gives
\(\eta+r\sim1\) and unit \(1-r\); unit \(r\) then gives
JO and infinitesimal \(r\) gives HO.
RO uses the IC parent dictionary with its newly proved
native bounds, not the narrower IC1 source premise.
JO and HO use \(j>0\), including \(j_0=0\) or slightly
negative.  Their exact three-circuit body and SB remote
proof retain all normal terms.

This proves the entire \(L,\eta,\zeta,j\) source
partition directly from LC4.  Every local source
estimate just invoked uses these native equations,
the original \(D,\beta\geq0\), its printed scale cell,
the NC gap and, where required, the unchanged
infinitesimal total scale.  No bound is obtained from
failure of an unnamed earlier theorem.
The seven restored exits are exactly large-\(L\)/unit-\(\eta\),
unit-\(L\)/small-\(\eta\), unit-\(L\)/unit-\(\eta\),
JR/JB, RO, JO and HO.  They are positive parts of this
cover.  Thus the actual JC constructions prove the
bare third row of LC5; any original earlier-complement
predicate can afterwards be intersected with the result.
The strengthening preserves physicality, signs, gaps,
literal boundaries and total scale.

All ten rows of LC5 now have complete original-parent
certificates.  They include every residual C1 source after
the positive-total-scale alternative, so C1 in LC1 is proved.

\subsection{Every C2 source and the absence of a circular premise}

The theorem CT1--CT16 proves the entire native C2 family.
Its proof takes the same total-scale split, treats both
weak complementary events by actual radial source entry,
and otherwise gives all nine ordered actual-angle types
in its explicit table.  Every empty cell has an exact
necessary-source contradiction; every occupied cell has
its complete local and remote full-parent proof.
In particular its directed weak HH cell is reflected to
the exact HJ domain, with \(C\) strong, \(\alpha\) weak
and both actual angles infinitesimal.
The HJ proof includes all ten source rows and every
allowed \(n=0\) map.  Its C1 boundary calls use only the
two explicit centered-A and small-\(e\) bare-source bodies,
never the C1 conclusion of the present theorem.
Thus CT has no dependence on this global C1/C2 join.
This proves C2 in LC1.

The mathematical dependencies have the following direction:
\[
\begin{gathered}
 \text{actual physical source geometry}
       \ \Longrightarrow\ {\rm SE,\ EN},\\
 \text{local identities and strict tensor Gram}
       \ \Longrightarrow\ {\rm AU,\ JC,\ HJ,\ UJ},
       \ \text{scoped terminals},\\
 \text{these terminals and finite source partitions}
       \ \Longrightarrow\ {\rm CT,\ LC}.
\end{gathered}\tag{LC8}
\]
A6 is separately proved from the actual zero-\(d\) geometry,
one- and two-hard-packet cases, weighted matching, promotion
and induction bodies; it does not use LC, CT, JC or HJ.
EN does not use AU or any tensor-SOS conclusion.
The source entries and coefficient-field physicality arguments
inside local proofs are not converted into positivity
statements in the tensor ring.
At no step is a mixed Schur denominator cancelled.
All roots and all first and higher state-normal terms belong
to the actual full parent LC2.

The finite mathematical inventory accompanying this proof
records direct and transitive actual proof bodies with their
restricted imported statements.  An old open-complement
paragraph, a historical status, a filename, or a checker
assertion count is not a mathematical premise.
Finite typed-source checks corroborate the exhaustion of
LC5 and CT6; strictness, conductor membership and full-normal
restoration are proved in the displayed mathematical bodies.

\subsection{Exact reconciliation of the historical rows}

The original thirteen-row checklist is reconciled by the following
restriction maps, not by changing its acceptance criterion.
Every row retains its original physical source, strict gaps,
infinitesimal total scale, testing/state domain and any originally
intersected earlier-complement predicate.
The present theorem is proved on the indicated bare domains,
so it also holds on those intersections.
The letter H in these original labels means an unbounded actual
half-angle, the type O used above.
\[
\begin{array}{l|l|l}
\text{original row}&\text{original role}&\text{present proof}\\ \hline
{\rm C1.1}&\text{earlier terminal}&{\rm LC5:\ C1.1\ exact\ bodies}\\
{\rm C1.2}&\text{earlier terminal}&{\rm LC5:\ centered\text{-}A\ theorem}\\
{\rm C1.3}&\text{earlier terminal}&{\rm LC5:\ small\text{-}e\ theorem}\\
{\rm C1.4}&\text{six refined rows}&{\rm JC\ and\ LC5\ bare\ scope}\\
{\rm C1.5}&\text{earlier terminal}&{\rm LC6\text{--}LC7,\ UI,\ UH}\\
{\rm C2.W.IH}&\text{earlier terminal}&{\rm CT6:\ weak\ IH\ full\ join}\\
{\rm C2.W.HI}&\text{empty source}&{\rm CT:\ WE1\text{--}WE6}\\
{\rm C2.W.HU}&\text{empty source}&{\rm CT:\ WE1\text{--}WE6}\\
{\rm C2.W.HH}&\text{ten refined rows}&{\rm CT4,\ HJ1\text{--}HJ16}\\
{\rm C2.S.II\_HH}&\text{earlier terminal}
                        &{\rm CT14:\ II\ local,\ remote,\ CT4}\\
{\rm C2.S.IH}&\text{earlier terminal}&{\rm CT:\ strong\ IH\ full\ join}\\
{\rm C2.S.UI\_HU}&\text{earlier terminal}
                        &{\rm CT:\ UI\ full\ join,\ CT4}\\
{\rm C2.S.HI}&\text{earlier terminal}&{\rm CT15:\ HI\ full\ join}
\end{array}\tag{LCH1}
\]
This explicitly accounts for the original nine earlier terminals,
the two empty source rows, and the two rows requiring subdivision.
The extra original condition \emph{at least one of \(m,A,D\) weak}
on each C2.S row is retained by restriction of the broader
both-strong \(C,\alpha\) statements in CT.  Neither emptiness
proof is replaced by an all-tests claim for nonexistent sources.
Every one of the nine earlier terminals has been reconnected
to its complete mathematical body and all necessary subordinate
proofs; its former status is not used.

For the C1.4 subdivision put \(L=k^2\zeta\), and, when
\(\eta\) is unbounded, \(\sigma=1/\eta\).
After the seven explicitly restored JC prior exits, the
original six predicates are exactly
\[
\begin{array}{l|l|l}
\text{label}&\text{retained source predicate}&\text{present theorem}\\ \hline
{\rm C14.1a}&L,\eta:{\rm O},\ h/\sigma\text{ bounded},
                  \ \beta\theta/\sigma^2:{\rm O}&{\rm JC}\\
{\rm C14.1b}&L,\eta:{\rm O},\ h/\sigma:{\rm O}&{\rm JC}\\
{\rm C14.2}&L:{\rm U},\ \eta:{\rm O}&{\rm JC}\\
{\rm C14.3a}&L:{\rm I},\ \eta\text{ bounded},\ \zeta:{\rm O}&{\rm JC}\\
{\rm C14.3b}&L,\eta:{\rm I},\ \zeta:{\rm U}&{\rm JC}\\
{\rm C14.4}&L:{\rm I},\ \eta:{\rm O},
                         \ \zeta\text{ noninfinitesimal}&{\rm JC}.
\end{array}\tag{LCH2}
\]
All six share the original C1.4 source input.  JC supplies
their separate complete local and outer constructions as
well as the precise scale disjunction, including its empty
intersection \(L\) unbounded with \(\eta\) infinitesimal.
The seven earlier source exits are proved, rather than
silently subtracted.  LCH2 is therefore an implication for
each original predicate, including its allowed literal
zero-event and source-opening faces.

For the weak-C2 subdivision use exactly the common II
coordinates HJ4--HJ10 after CT4:
\[
 z=u_c^2,\quad y=v_c^2,\quad B=H+L,\quad
 a=z/B,\quad t=y/B,\quad h=H/B,\quad l=L/B,\quad e=1-r.
\]
Here \(B,z,y\) are positive infinitesimal,
\(h+l=r+e=1\), all original gaps remain positive and
the actual \(n\geq0\) formula is retained.  The ten
original predicates and present proofs are
\[
\begin{array}{l|l|l}
{\rm R1}&t:{\rm I}&{\rm HJ}\\
{\rm R2}&t:{\rm O}&{\rm HJ}\\
{\rm R3}&t,h:{\rm U},\ l:{\rm I}&{\rm HJ}\\
{\rm R4}&t,h,l:{\rm U},\ j:{\rm I}&{\rm HJ}\\
{\rm R5}&t,h,l,j:{\rm U},\ r:{\rm I}&{\rm HJ,\ BSR}\\
{\rm R6}&t,h,l,j,r:{\rm U},\ g:{\rm I}&{\rm HJ,\ BSR}\\
{\rm R7}&t:{\rm U},\ h,g:{\rm I},\
                     r,j\text{ not both I}&{\rm HJ,\ BSR}\\
{\rm R8}&t,g:{\rm U},\ h,j:{\rm I}&{\rm HJ,\ BSR}\\
{\rm R9}&t,g,j:{\rm U},\ h,e,a:{\rm I}&{\rm HJ,\ BSR}\\
{\rm R10}&t,g,j:{\rm U},\ h,e:{\rm I},\ a:{\rm O}&{\rm HJ,\ BSR}.
\end{array}\tag{LCH3}
\]
These rows use the original ordered dispatch, not an
unstated quotient of the directed angle list.
HJ14 proves that the all-six-unit alternative contradicts
the retained infinitesimal total scale.
HJ15 proves that \(r,g,h,j\) all infinitesimal would make
\(C\) weak, contradicting this directed strong-\(C\)
input; it justifies the exclusion written in R7.
The other earlier exits are dispatched by their actual
source formulas and complete full-parent bodies in HJ/BSR.
Thus the ten-row exhaustion and each terminal implication,
not merely the number ten, are supplied.

\subsection{Restoration and the universal slack consequence}

Every local proof gives the unmultiplied parent, or an
equivalent parent multiplied only by a positive pure-source
or pure-testing factor whose nonvanishing is explicitly
proved.  Such a factor is removed by its square root inside
that field.  Original retained packets are then added with
their actual nonnegative coefficients.
The exact homogeneous angle identities restore every finite
chart and include literal endpoints without division by
a vanishing coordinate.  Pure-testing state normalizations
cover all nonzero real projective states, and the zero state
is immediate.  Since LC2 is a real quadratic form, applying
the real result to real and imaginary components and adding
proves the complex-state assertion.  Restoring \(d\)
finishes the complete local-certificate theorem.

\begin{corollary}[Universal separated slack]
Let \(\mathcal Q_{222}\subset\mathbb R^8\) be the finite-dimensional
bipartite tensor-product quantum set in the two-setting,
two-outcome scenario, and let \(P\) be its cone of nonnegative
affine functionals.  Use the fixed finite real-qubit formula
for \((\mathcal Q_{222})_{K_t}\) and the transferred cone
\(P_{K_c}\).  Then for every separately labelled \(K_c,K_t\)
as above,
\[
 f_0\otimes1+\sum_{\nu=1}^8 f_\nu\otimes q_\nu
       \ \in\ \Sigma T^2
 \qquad
 (f\in P_{K_c},\ q\in(\mathcal Q_{222})_{K_t}).
\tag{LC9}
\]
\end{corollary}

\begin{proof}
The complete source reduction SE11--SE15 writes every normalized
source in \(K_c\) as a finite conic combination of event atoms,
the two explicit noncoercive closed faces, and the actual
canonical contact sources LC1.  It proves this fixed finite
statement over the reals and transfers it; no arbitrary
infinite extreme-point expansion is used.
Event atoms have separated squares.  Each of the two closed
faces has the actual finite matrix lift and strict feasible
point in \emph{closed\_faces\_short.tex}.
Attained finite-dimensional duality and separate-field
positive-semidefinite factorizations give its tensor
trace-pairing squares.  Canonical atoms use the theorem
just proved, followed by the exact homogeneous slack
normalization in the physical reduction.

Each testing point is a convex combination of at most nine
real-qubit atoms in its own field.  If the source combination
weights are \(a_j\geq0\), testing weights are \(w_i\geq0\)
with \(\sum_iw_i=1\), and each atom pairing is
\(\sum_k h_{ijk}^2\), then the full pairing is exactly
\[
       \sum_{i,j,k}
        \left[(\sqrt{a_j}\otimes\sqrt{w_i})h_{ijk}\right]^2 .
\tag{LC10}
\]
Affineness preserves the constant term because the testing
weights sum to one.  Zero weights are never inverted.
All Bell relabelings and all positive normalizing factors
are restored inside their own field.  This is finite conic
and convex closure in the tensor ring, and proves LC9.
\end{proof}

The proof supplies the universal tensor statement, not an
evaluated affine matrix pencil or a numerical bound on its
size.  The separate finite-lift criterion applies to LC9
using two labelled copies of any one real closed extension;
its compactness conclusion is not inferred from a sum of
the local root counts.
